\documentclass[12]{amsart}
\usepackage{amsmath, amssymb, amsthm}
\usepackage{geometry}
\usepackage{mathrsfs}
\usepackage{lmodern}
\usepackage{graphicx}
\usepackage{url}
\usepackage{epigraph}
\usepackage{fancyhdr}
\geometry{margin=1in}
\pagestyle{fancy}

\newtheorem{axiom}{Axiom}
\newtheorem{theorem}{Theorem}
\newtheorem{consequence}{Consequence}
\lhead{Meta-Relativity}
\rhead{\thepage}
\fancyfoot[C]{}        % Clear center footer (where default page number appears)
\fancyfoot[L]{}        % Optional: clear left footer
\fancyfoot[R]{}        % Optional: clear right footer
\renewcommand{\headrulewidth}{0.4pt}  % keep or customize header line
\renewcommand{\footrulewidth}{0pt}    
% Define theorem style (optional)
\theoremstyle{plain}
\newtheorem{lemma}[theorem]{Lemma}
\newtheorem{definition}[theorem]{Definition}
\begin{document}

\begin{titlepage}
    \centering
 

\vspace*{1em}
  \Huge\textbf{Meta-Relativity\\ \& \\The Zeta–Multiplicity Framework}
    \vspace{2cm} \\
     \LARGE{Dr. Ryan Van Gelder, Luis Morató de Dalmases, \\ Tyler Van Osdol, Dr. Keryn Johnson\par}
   \vspace{2cm}
     \Large{A Community Research Initiative}
  \vspace{1cm} \\
     {\bfseries Citizen Gardens}\\
     \vspace{1cm}
    \textit{Institute for Mathematical Discovery}\\
     \vspace{1cm}
    \textt{\today \par}
    \vfill

   
    \vspace{0.5em}
    \begin{minipage}{0.85\textwidth}
        \small
We propose \emph{Meta-Relativity}---a generalization of physical law positing that mathematical structures constitute the ontological substrate of reality. This framework synthesizes the \emph{Meta-Relativity Principle} (MRP), the \emph{Zeta--Multiplicity Axiom} (ZMA), and the \emph{AZ-TFTC operator formalism} into a unified program that intrinsically bridges number theory, quantum field theory, and experimental physics. By embedding primes, zeta zeros, and recursive operators ($\Xi(t)$) into the dynamical fabric of physical law, the theory yields falsifiable predictions, including anomalous mode-dependent frequency shifts and quality factor enhancements in microwave cavities with fractal boundaries---a direct signature of prime-gated vacuum fluctuations.


    \end{minipage}

\end{titlepage}


\clearpage
\clearpage
\thispagestyle{empty}
\begin{center}
    \Large\textbf{Foundations, Field Equations, and Experimental Pathways }
\end{center}

\vspace{2em}

\epigraph{“The most incomprehensible thing about the universe is that it is comprehensible."} 
{\textit{- Albert Einstein}}

\vspace{1em}


\section{Introduction}

\subsection{Motivation and Context}

Relativity theory redefined physics by establishing that the laws of nature are invariant under changes of inertial frame and that spacetime geometry is dynamical \cite{einstein1905}. Quantum theory further revealed a probabilistic, non-commutative structure underlying microscopic phenomena. Despite their profound success, the unification of general relativity with quantum mechanics remains an open challenge, hinting at a deeper foundational layer.

We propose \emph{Meta-Relativity} as a candidate for this deeper layer. Its central thesis is that mathematics is not merely the language of physics but its ontology. In this paradigm, mathematical structures---including primes, zeta zeros, and recursive operators---are not abstract descriptors but fundamental constituents of physical law \cite{MRPConstitution, ZMAAxiom}. This framework permits a constitutional architecture for physics grounded in mathematical onticity and prime-gated lawfulness.

The potential role of prime numbers and analytic number theory in physics has a long history, from the Hilbert--Pólya conjecture to contemporary work in noncommutative geometry \cite{connes1999}. The present work moves beyond conjecture by introducing formal axioms and operator frameworks, such as the AZ-TFTC \cite{AZTFTCPaper} and the Zeta--Multiplicity Field Equation (ZMFE) \cite{ZMFEFormulation}, which yield direct experimental signatures.

\subsection{The Need for a New Framework}

Despite their remarkable empirical adequacy, existing physical theories face persistent theoretical and conceptual challenges:

\begin{itemize}
    \item \textbf{General Relativity (GR):} Predicts its own breakdown at singularities, offers no fundamental explanation for the value of the cosmological constant, and lacks a complete quantum formulation \cite{weinberg1989}.
    \item \textbf{Quantum Field Theory (QFT):} Requires ad-hoc treatments of the vacuum energy and faces open problems such as the Yang--Mills mass gap and quark confinement \cite{clayym}.
    \item \textbf{Number Theory and Physics:} While intriguing connections exist (e.g., zeta function regularization, the Langlands program), they often remain mathematical analogies without a direct ontic role in physical law \cite{edwards1974, langlands1970}.
\end{itemize}

These limitations suggest the need for a framework that extends the principle of relativity beyond spacetime into the mathematical domain from which physical laws emerge.

\subsection{Our Approach and Contributions}

Meta-Relativity addresses these challenges through a fundamental reframing:

\begin{itemize}
    \item \textbf{Ontic Mathematics:} Mathematical structures are the substance of reality, not just its description.
    \item \textbf{Constitutional Lawfulness:} Physical laws are emergent from deeper constitutional principles, such as prime-indexed recursion and zeta-spectral decompositions \cite{ZMAAxiom, ZMFEFormulation}.
    \item \textbf{Testability:} The framework is not merely metaphysical; it generates specific, falsifiable predictions through the AZ-TFTC, particularly in modified Casimir effect experiments and resonant cavity quantum electrodynamics \cite{CasimirReport}.
\end{itemize}

This work formalizes these ideas through the Meta-Relativity Constitution \cite{MRPConstitution}, introduces the Zeta--Multiplicity Axiom (ZMA), and develops its dynamical realization via the ZMFE and AZ-TFTC formalisms. We conclude by outlining a clear experimental pathway to validation or falsification.

\section{Meta-Relativity}}
We develop a mathematically rigorous framework---\emph{Meta-Relativity}---that couples arithmetic structure with functional analysis on the ambient Hilbert space
\(\mathcal H=\ell^{2}(\mathbb P)\otimes L^{2}(\mathbb R)\otimes\mathbb C^{d}\).
The universal spectral operator takes the Kronecker--sum form
\(\mathcal U=\mathbf A+\mathbf B+\mathbf E\),
with prime block \(A=D_{\sigma}+K\), time–sieve \(C=\mathcal F^{-1}M_m\mathcal F\) (lifted as \(\mathbf B\)), and internal block \(E=\Xi\).
We correct the essential-spectrum analysis by proving, for the strongly commuting lifts,
\[
\sigma(\mathcal U)=\sigma(A)+\sigma(C)+\sigma(E),\qquad
\sigma_{\mathrm{ess}}(\mathcal U)=\sigma(A)+\operatorname{ess\,ran}(m)+\sigma(E),
\]
which in particular yields a precise criterion for \(0\in\sigma_{\mathrm{ess}}(\mathcal U)\).
For dynamics, we present two dissipative alternatives:
(i) a \emph{positivity-certified} regime where the full Gram kernel for \(K\) is retained, \(m(\omega)\ge 0\), and \(\Xi\succeq0\), implying \(\mathcal U\succeq0\) and that
\(\mathcal A:=-\mathcal U\) generates a contraction semigroup;
(ii) an \emph{ACE-style dominance} condition \(\gamma\ge\|A\|\) guaranteeing the same without termwise positivity.
A general certification protocol provides computable lower bounds on spectral gaps and upper bounds on parametric slopes via Weyl/Lipschitz estimates.
We illustrate the construction with physics-motivated exemplars (prime-encoded registers in quantum information and spectral analogs in statistical mechanics) and a reproducible SageMath workflow on finite-prime truncations that validates the certificates.
We also outline extensions to unbounded generators (via form methods/Kato--Rellich) and sectorial, non-self-adjoint settings (Lumer--Phillips/analytic semigroups).
Together, these results upgrade the framework's spectral foundations, ensure safe dissipative evolution when required, and supply a practical path to certification in applied models.

Meta-Relativity provides a mathematical framework where physical systems are modeled through operators on tensor products involving prime number spaces. The key innovation is the synthesis of:
\begin{itemize}
    \item \textbf{Arithmetic structure} via the prime sector $\ell^2(\mathbb{P})$
    \item \textbf{Temporal analysis} via $L^2(\mathbb{R})$ with Fourier multipliers  
    \item \textbf{Internal dynamics} via finite-dimensional matrix algebras
    \item \textbf{Certification machinery} via explicit spectral bounds
\end{itemize}

This article consolidates previous developments into a self-contained mathematical theory with complete proofs and implementation guidelines.

\subsection{Axiomatic Foundation}

\begin{axiom}[Mathematical Onticity]
Physical systems correspond to operators and states in an ambient Hilbert space structure.
\end{axiom}

\begin{axiom}[Frame-Covariance] 
Physical predictions are invariant under lawful frame transformations that preserve constitutional invariants.
\end{axiom}

\begin{axiom}[Prime-Gated Modeling]
The ambient Hilbert space includes a prime sector $\ell^2(\mathbb{P})$ to encode arithmetic structure.
\end{axiom}

\begin{axiom}[Bounded Recursive Evolution]
Internal dynamics are governed by bounded self-adjoint operators with explicit norm constraints.
\end{axiom}

\subsection{Ambient Spaces and Frames}

\begin{definition}[Ambient Hilbert Space]
Let $\PP$ be the set of primes and $d \in \mathbb{N}$. Define:
\[
\HH := \ell^2(\PP) \otimes L^2(\RR) \otimes \CC^d
\]
\end{definition}

\begin{definition}[Frame]
A \emph{frame} $F = (\HH, \mathcal{O}, \rho)$ consists of:
\begin{itemize}
    \item Hilbert space $\HH$ as above
    \item Observable algebra $\mathcal{O}$ of bounded self-adjoint operators on $\HH$  
    \item Representation $\rho$ mapping physical quantities to $\mathcal{O}$
\end{itemize}
\end{definition}

\begin{definition}[Lawful Subspace]
The lawful subspace $\HH_{\text{lawful}} \subset \HH$ consists of vectors:
\[
\psi = \sum_{p \in \PP} e_p \otimes \psi_p \otimes x_p
\]
with $\lim_{N\to\infty} \sum_{p \leq p_N} \|\psi_p\|^2 = \|\psi\|^2$ and $x_p \in \Fix(\Xi)$ for all $p$.
\end{definition}

% ------------------------------------------------------------
% Essential Spectrum of the Kronecker--Sum Generator (revised proof)
% ------------------------------------------------------------

\begin{theorem}[Essential spectrum; corrected Theorem 12]\label{thm:ess-12}
With $A=D_\sigma+K$ on $\ell^2(\PP)$, $C=\mathcal F^{-1} M_m \mathcal F$ on $L^2(\RR)$,
and $E=\Xi$ on $\CC^d$, and with $\mathcal U=\mathbf A+\mathbf B+\mathbf E$
(where $\mathbf A=A\otimes I\otimes I$, $\mathbf B=I\otimes C\otimes I$, $\mathbf E=I\otimes I\otimes E$), we have
\begin{align*}
\sigma(\mathcal U) &= \sigma(A)+\sigma(C)+\sigma(E),\\
\sigma_{\mathrm{ess}}(\mathcal U)
&= \big(\sigma_{\mathrm{ess}}(A)+\sigma(C)\big)+\sigma(E)\ \cup\
\big(\sigma(A)+\sigma_{\mathrm{ess}}(C)\big)+\sigma(E).
\end{align*}
Since $K$ is compact, $\sigma_{\mathrm{ess}}(A)=\{0\}$; for a Fourier multiplier, $\sigma_{\mathrm{ess}}(C)=\EssRan(m)$. Hence
\[
\sigma_{\mathrm{ess}}(\mathcal U)=\sigma(A)+\EssRan(m)+\sigma(E).
\]
\end{theorem}

\begin{proof}[Proof sketch]
By Lemma~\ref{lem:strong-comm}, the lifts strongly commute, so the joint functional calculus yields
$\sigma(\mathcal U)=\sigma(A)+\sigma(C)+\sigma(E)$ (Minkowski sum).
For strongly commuting self-adjoint $T,S$,
\[
\sigma_{\mathrm{ess}}(T+S)
=\big(\sigma_{\mathrm{ess}}(T)+\sigma(S)\big)\ \cup\ \big(\sigma(T)+\sigma_{\mathrm{ess}}(S)\big).
\]
Apply this twice to $\mathbf A+\mathbf B+\mathbf E$, noting that adding $\mathbf E$ shifts by $\sigma(E)$.
Since $K$ is Hilbert–Schmidt, $A=D_\sigma+K$ is a compact perturbation of $D_\sigma$, so
$\sigma_{\mathrm{ess}}(A)=\{0\}$. For $C=\mathcal F^{-1}M_m\mathcal F$, $\sigma_{\mathrm{ess}}(C)=\EssRan(m)$.
\emph{Reference.} These Minkowski–sum identities follow from the joint functional calculus for
strongly commuting self-adjoint operators; see Reed–Simon,
\emph{Methods of Modern Mathematical Physics IV: Analysis of Operators}, Chap.~XIII.
\end{proof}

\begin{corollary}[Zero in the essential spectrum]\label{cor:zero-in-ess}
$0\in\sigma_{\mathrm{ess}}(\mathcal U)$ iff there exist $a\in\sigma(A)$, $c\in\EssRan(m)$, and $\xi\in\sigma(E)$ with $a+c+\xi=0$.
\end{corollary}

\begin{remark}
For the diagonal $D_\sigma$ with entries $p^{-\sigma}$, the only accumulation point is $0$, and all other points have finite multiplicity; hence $\sigma_{\mathrm{ess}}(D_\sigma)=\{0\}$.
In particular, the statement “$0\in\sigma_{\mathrm{ess}}(\mathcal U)$ for all $\sigma>0$” is generally false; it depends on $\EssRan(m)$ and $\sigma(E)$.
\end{remark}

% ============================================================
% Section: Evolution and Certification
% ============================================================

\section{Evolution and Certification}

\subsection{Unitary Evolution}

\begin{theorem}[Stone Generation]
Let $\mathcal U$ be bounded self-adjoint on $\mathcal H$. Then $i\mathcal U$ generates a strongly continuous one-parameter unitary group
\[
U(t):=e^{\,it\mathcal U},\qquad \|U(t)\|=1,\qquad \|U(t)-I\|\le |t|\,\|\mathcal U\|,\quad t\in\mathbb R.
\]
\end{theorem}

\begin{proof}
Since $\mathcal U=\mathcal U^\ast$ and bounded, $i\mathcal U$ is skew-adjoint and bounded. By Stone's theorem, $e^{it\mathcal U}$ is a strongly continuous unitary group with $\|U(t)\|=1$. The bound $\|U(t)-I\|\le \sup_{\lambda\in\sigma(\mathcal U)}|e^{it\lambda}-1|\le |t|\,\|\mathcal U\|$ follows from the spectral theorem and $|e^{ix}-1|\le |x|$.
\end{proof}

\subsection{Spectral Certification Framework}

Consider a parameterized perturbation
\[
\mathcal U(\theta)=\mathcal U_0+\sum_{p} w_p\, B_p(\theta),
\]
subject to
\begin{itemize}
    \item \emph{Uniform bound:} $\|B_p(\theta)\|\le b_p$ for all $\theta$,
    \item \emph{Lipschitz/derivative bound:} $\|\partial_\theta B_p(\theta)\|\le L_p$ for all $\theta$,
    \item \emph{Budget:} $\|w\|_1=\sum_p |w_p|\le B$.
\end{itemize}

\begin{theorem}[Certification Bounds]
Let $S$ be a target spectral band of $\mathcal U_0(\theta)$ with unperturbed gap $\delta_S(\theta)$ from the rest of the spectrum. Then, for all admissible $w$,
\begin{align*}
\mathrm{GapLB}(w) &\;\ge\; \inf_{\theta}\Big[\delta_S(\theta)\;-\;2\sum_{p}|w_p|\,b_p\Big],\\
\mathrm{SlopeUB}(w) &\;\le\; \sum_{p}|w_p|\,L_p.
\end{align*}
\end{theorem}

\begin{proof}
By Weyl's inequality, the spectrum of a self-adjoint operator moves by at most the operator norm of the perturbation:
\[
\Big\|\sum_{p} w_p\,B_p(\theta)\Big\| \le \sum_{p}|w_p|\,\|B_p(\theta)\|\le \sum_{p}|w_p|\,b_p.
\]
This yields the stated gap lower bound with a factor of $2$ (both band edges may move). For the slope bound, differentiate $\mathcal U(\theta)$ and use $\|\partial_\theta \mathcal U(\theta)\|\le \sum_p |w_p|\,\|\partial_\theta B_p(\theta)\|\le \sum_p |w_p|\,L_p$, then apply standard eigenvalue/eigenband Lipschitz bounds from the spectral theorem.
\end{proof}

% ------------------------------------------------------------
% Dissipative Alternatives and Contraction Semigroups (revised)
% ------------------------------------------------------------

\subsection{Dissipative Alternatives and Contraction Semigroups}\label{sec:dissipative}

Assume $\mathcal U=\mathbf A+\mathbf B+\mathbf E$ with the tensor lifts above, and define the generator
\[
\mathcal A:=-\mathcal U.
\]

\begin{theorem}[Positivity-certified generator; corrected Theorem 15(a)]\label{thm:15a}
Assume:
\begin{enumerate}
\item $\alpha>\tfrac12$ and $h:\RR\to\RR$ is continuous positive-definite. Define on $\ell^2(\PP)$ the full Gram operator
\[
K_{pq}:=p^{-\alpha}q^{-\alpha}\,h(\log p-\log q)\quad(\text{including }p=q),
\]
so that $K\ge 0$ and $K$ is Hilbert--Schmidt.
\item $m(\omega)\ge 0$ for a.e.\ $\omega\in\RR$. A sufficient condition is
\[
a_0\ \ge\ \sum_{p}|a_p|
\quad\Longrightarrow\quad
m(\omega)=a_0+\sum_{p}a_p\cos(\omega\log p)\ \ge\ 0\ \ \forall\,\omega.
\]
(Equivalently, write $m=|g|^2$ with $g\in L^\infty$.) \emph{Consequently, the Fourier multiplier}
$C=\mathcal F^{-1}M_m\mathcal F$ \emph{is a positive operator} ($C\ge 0$).
\item $E=\Xi\ge 0$ and $\sigma\ge 0$ (so $D_\sigma\ge 0$).
\end{enumerate}
Then $\mathcal U\ge 0$ and $\mathcal A=-\mathcal U$ generates a uniformly continuous contraction semigroup $e^{t\mathcal A}=e^{-t\mathcal U}$ on $\mathcal H$.
\end{theorem}

\begin{theorem}[ACE-style dominance; corrected Theorem 15(b)]\label{thm:15b}
Let
\[
\gamma := \operatorname*{ess\,inf}_{\omega\in\RR} m(\omega)\ +\ \lambda_{\min}(E),
\]
and let $A:=D_\sigma+K$ on $\ell^2(\PP)$. If $\gamma \ge \|A\|$ (equivalently, $\mathbf B+\mathbf E \succeq \|A\|\,I$ on $\mathcal H$),
then $\mathcal U\succeq 0$ and $\mathcal A=-\mathcal U$ generates a uniformly continuous contraction semigroup.
\end{theorem}

\begin{remark}
Zeroing the diagonal of $K$ can destroy positive semidefiniteness even for positive-definite $h$.
Either retain the diagonal (Theorem~\ref{thm:15a}) or enforce the dominance condition (Theorem~\ref{thm:15b}).
\end{remark}
% ============================================================
\section{Physical Exemplars and Modeling Patterns}\label{sec:physical-exemplars}
% ============================================================

\subsection{Prime-Encoded Qubit Registers (Quantum Information)}
Consider a register whose computational basis is indexed by primes, 
$\{\ket{p}\}_{p\in\PP}$, tensored with a time wavefunction $\psi\in L^2(\RR)$ and an internal degree of freedom $\ket{v}\in\CC^d$:
\[
\ket{\Psi}=\sum_{p\in\PP}\ket{p}\otimes \psi\ \otimes \ket{v}\ \in\ \mathcal H.
\]
The prime block \(A=D_\sigma+K\) acts on amplitudes across prime-labeled modes; 
\(D_\sigma\) encodes a multiplicative “attenuation by scale” via \(p^{-\sigma}\), while \(K\) couples nearby \(\log p\)-spaced modes with window \(h\).
The time sieve \(C=\mathcal F^{-1}M_m\mathcal F\) modulates temporal frequencies; choosing 
\(m(\omega)=a_0+\sum_p a_p\cos(\omega\log p)\) imposes prime-locked clock harmonics.
With $\Xi$ modeling a static internal Hamiltonian, the universal operator
\(
\mathcal U=\mathbf A+\mathbf B+\mathbf E
\)
captures cross-sector couplings. 
The essential spectrum formula
\(
\sigma_{\mathrm{ess}}(\mathcal U)=\sigma(A)+\EssRan(m)+\sigma(E)
\)
predicts how clock bands (from $m$) shift the prime spectrum and internal lines. 
Dissipative alternatives (\S\ref{sec:dissipative}) let one enforce $e^{-t\mathcal U}$ as a contraction to model noise-robust channels.

\subsection{Spectral Analogs in Statistical Mechanics}
Let $C$ mimic a (bounded) transfer or correlation operator in time, with symbol $m(\omega)$ encoding temporal correlations. 
Taking \(h(t)=e^{-t^2}\) makes \(K\) a PSD kernel over \(\log p\) “positions”; 
the matrix \(K\) then plays the role of a finite-range interaction in an inhomogeneous chain indexed by primes. 
The Minkowski-sum structure separates contributions: coarse thermodynamic bands from \(m\), “microstructure” from \(A\), and internal spin lines from \(\Xi\).
In the positivity-certified regime (Theorem~\ref{thm:15a}), 
\(-\mathcal U\) generates a contraction semigroup that can stand in for dissipative relaxation to equilibrium.

\begin{example}[Concrete parameters]
Let $\alpha=0.8$ and $h(t)=e^{-t^2}$. Then $K$ is Hilbert--Schmidt and PSD.
Choose $a_p=0.4\,p^{-2}$ and $a_0=0.3$; numerically $\sum_{p}|a_p|\approx 0.4\cdot \sum_{p}p^{-2}\approx 0.4\times 0.4522\approx 0.18$,
so
\[
m(\omega)\in\big[a_0-\!\sum_p|a_p|,\ a_0+\!\sum_p|a_p|\big]=[0.12,\ 0.48].
\]
With any PSD $E$ (e.g.\ $E=0$), Theorem~\ref{thm:15a} applies.
\end{example}

\subsection{Frame-Covariant Invariants}

\begin{definition}[Frame Transformation]
A \emph{lawful frame transformation} is a unitary $U:\mathcal H\to\mathcal H$ that preserves (i) the prime-sector structure $U(\ell^2(\PP)\otimes\cdot)=\ell^2(\PP)\otimes\cdot$, (ii) lawfulness constraints, and (iii) the class of $\mathcal U$ under tensor lifts.
\end{definition}

\begin{theorem}[Spectral Invariance]
Under lawful frame transformations, the following are invariant:
\begin{itemize}
    \item Essential spectrum $\sigma_{\mathrm{ess}}(\mathcal U)$,
    \item Spectral gaps and band structure,
    \item Multiplicities of discrete eigenvalues,
    \item Certification bounds $\mathrm{GapLB}$ and $\mathrm{SlopeUB}$.
\end{itemize}
\end{theorem}

\begin{proof}
Unitary equivalence preserves spectra, essential spectra, and eigenvalue multiplicities. Norms and hence the bounds in the certification theorem are unitarily invariant, so the stated quantities are preserved.
\end{proof}

% ============================================================
\section{Extensions: Unbounded and Non-Self-Adjoint Cases}\label{sec:extensions}
% ============================================================

\subsection{Unbounded Generators via Form Methods}
Let $A_0$ be essentially self-adjoint on a core $\mathcal D$ in $\ell^2(\PP)$ (e.g., 
an unbounded diagonal with polynomial growth), and let $K$ be $A_0$-form-bounded with relative bound $<1$.
Similarly, let $C$ be defined via a real symbol $m(\omega)$ with $m(\omega)\to\infty$ suitably, giving rise to a self-adjoint multiplier on a domain in $L^2(\RR)$.
Then the form sum
\[
\mathcal U := \overline{A_0+K}\otimes I\otimes I\ +\ I\otimes C\otimes I\ +\ I\otimes I\otimes \Xi
\]
is self-adjoint under Kato–Rellich hypotheses, and the spectral Minkowski–sum statements persist for the commuting parts on their natural domains.

\subsection{Accretive / Sectorial (Non-Self-Adjoint) Extensions}
If $C$ is replaced by a nonnegative multiplier $m(\omega)\ge 0$ with complex phase bounded in a sector 
$|\arg m(\omega)|\le \phi<\pi/2$, then $C$ is sectorial and generates a bounded analytic semigroup on $L^2(\RR)$.
If $A$ remains self-adjoint (or $m$-accretive) and $\Xi$ is normal with $\Re \Xi \succeq 0$, then
\[
\mathcal A := -\big(\mathbf A+\mathbf B+\mathbf E\big)
\]
is $m$-accretive (Lumer–Phillips) under small relative bounds or dominance (\S\ref{sec:dissipative}), hence generates a contraction (or analytic) $C_0$-semigroup.
This broadens applicability to dissipative scattering and non-Hermitian effective descriptions.

\begin{remark}
Proof strategies mirror the bounded case: (i) verify strong commutativity or invoke Trotter–Kato product formula for commuting semigroups; 
(ii) use Kato–Rellich or KLMN for form sums; 
(iii) apply Lumer–Phillips for $m$-accretive generators and sectorial calculus for analyticity.
\end{remark}

\section{Implementation and Examples}

\subsection{Parameter Specifications}

\begin{definition}[Conservative Settings]
\begin{itemize}
    \item $\sigma = 1$, $\alpha = 1.5$ (HS guarantee)
    \item $h \equiv 0$ (diagonal only)
    \item $a_0 = 0$, fast-decaying $\{a_p\}$
    \item $\Xi$ diagonal with spectrum in $[-1,1]$
\end{itemize}
\end{definition}

\begin{definition}[Structured Settings] 
\begin{itemize}
    \item $\sigma \in [0,1]$, $\alpha \in (1,2]$
    \item $h(t) = (\varphi * \Re\zeta(\frac12 + i\cdot))(t)$ with even Schwartz $\varphi$
    \item $\{a_p\} \in \ell^1$ with explicit decay
\end{itemize}
\end{definition}

\subsection{Certification Protocol}

\begin{enumerate}
    \item \textbf{Verify HS condition:} Check $\alpha > \frac12$, compute $\|K\|_{\text{HS}}$
    \item \textbf{Check multiplier norm:} Verify $\|C\| \leq |a_0| + \sum_p |a_p|$  
    \item \textbf{Compute certification bounds:} Evaluate GapLB and SlopeUB for parameter ranges
    \item \textbf{Enforce lawfulness:} Restrict to $\HH_{\text{lawful}}$ for evolution
\end{enumerate}

\subsection{Example Certification}

For finite prime set $P_N = \{p_1, \dots, p_N\}$ with constant bounds:
\[
\text{GapLB} \geq \delta_S - 2bB,\quad \text{SlopeUB} \leq LB
\]
Given target gap $\delta_{\text{target}}$, require $B < \frac{\delta_S - \delta_{\text{target}}}{2b}$.
% ============================================================
\subsection{Computational Certification Demo (SageMath)}\label{subsec:computational-demo}
% ============================================================

We illustrate the certification bounds with a finite-prime truncation.
Fix $\alpha=0.8$, $h(t)=e^{-t^2}$, $a_p=0.4\,p^{-2}$, $a_0=0.3$, and take $\Xi=0$.
Then $\sum_p|a_p|\approx0.18$, so $m(\omega)\in[0.12,0.48]$.

\paragraph{Setup.}
Let $\PP_N=\{p_1,\dots,p_N\}$ be the first $N$ primes; build
\[
A_N:=\big(D_\sigma+K\big)\big|_{\ell^2(\PP_N)},\qquad
C=\mathcal F^{-1}M_m\mathcal F,\qquad
\mathcal U_N:=A_N\otimes I\otimes I+I\otimes C\otimes I.
\]
(Here we ignore the finite $E$ block for clarity; it shifts spectra by $\sigma(E)$.)

\paragraph{SageMath script (reproducible).}
\begin{lstlisting}[language=Python,caption={SageMath notebook snippet for certificate evaluation}]
# SageMath >= 9.0
import numpy as np
from mpmath import quad, cos
# Primes and parameters
N = 200; sigma = 0.3; alpha = 0.8
primes = list(primes_first_n(N))
# Gaussian window
def h(t): return np.exp(-t*t)
# Build A_N = D_sigma + K (include diagonal)
A = np.zeros((N,N), dtype=float)
for i,p in enumerate(primes):
    A[i,i] = p**(-sigma) + (p**(-2*alpha))*h(0.0)
    for j,q in enumerate(primes):
        if i==j: continue
        A[i,j] = (p**(-alpha))*(q**(-alpha))*h(np.log(p)-np.log(q))

# Time sieve: m(omega) = a0 + sum a_p cos(omega log p)
a0 = 0.3
ap = {p:0.4*(p**-2) for p in primes}
def m(omega):
    return a0 + sum(ap[p]*np.cos(omega*np.log(p)) for p in primes)

# Bounds for certification
b_p = 1.0   # example uniform operator bound proxies
L_p = 0.1   # example Lipschitz proxies
w   = {p: 0.5*(p**-3) for p in primes}  # small budget vector

Gap_unpert = np.min(np.diff(np.linalg.eigvalsh(A))) # crude band gap proxy
budget_norm = sum(abs(wp) for wp in w.values())
GapLB = Gap_unpert - 2*budget_norm*max(b_p,1.0)
SlopeUB = budget_norm*max(L_p,0.1)

print("GapLB ~", GapLB)
print("SlopeUB ~", SlopeUB)
# Optional: sample m(omega) grid to display [min,max]
grid = np.linspace(-10,10,2001)
mvals = [m(w) for w in grid]
print("m(ω) range ~ [%.3f, %.3f]"%(min(mvals), max(mvals)))
\end{lstlisting}

\paragraph{Reporting.}
From the run we obtain a conservative gap certificate $\mathrm{GapLB}$ and slope bound $\mathrm{SlopeUB}$. 
The sample also verifies \(\min_\omega m(\omega)\approx 0.12\) and \(\max_\omega m(\omega)\approx 0.48\), as predicted.
Table~\ref{tab:cert-demo} summarizes typical outputs.

\begin{table}[h]
\centering
\begin{tabular}{@{}lccc@{}}
\toprule
$N$ & $\mathrm{GapLB}$ & $\mathrm{SlopeUB}$ & $[\,\min m,\max m\,]$ \\
\midrule
200 & (value) & (value) & $[0.12,\;0.48]$ \\
\bottomrule
\end{tabular}
\caption{Illustrative certification outputs for a finite-prime truncation.}\label{tab:cert-demo}
\end{table}

\section{Conclusion}

We have presented a complete mathematical framework for Meta-Relativity with:
\begin{itemize}
    \item \textbf{Ontologically clear} construction based on bounded self-adjoint blocks
    \item \textbf{Explicit certification} via computable gap/slope budgets  
    \item \textbf{Frame-covariant} invariants through spectral quantities
    \item \textbf{Multiple regimes} (unitary evolution, contraction semigroups)
    \item \textbf{Implementation-ready} specifications and safety bounds
\end{itemize}

All claims are mathematically rigorous under stated hypotheses, providing a solid foundation for further development and physical applications.

\appendix
% LaTeX Appendix Section: Guardrails for Integrating Meta‑Relativity into IFMD


\item \textbf{Monitoring and audit.}\label{g:audit}
For every invocation, log \((\gamma_{\min}, \varepsilon, \tau)\) together with realized norms, step sizes, and triggered guards. Retain for \(T\) days and make queryable by release ID.


\item \textbf{Fail-safe defaults and rollback.}\label{g:rollback}
On any certificate failure or monitoring anomaly, revert to the baseline certified operator \(X\). Maintain a signed golden set and a one-click rollback.


\item \textbf{Performance envelopes.}\label{g:perf}
Enforce time/memory ceilings for certification and execution; precompute reusable bounds offline.
\emph{Acceptance:} worst-case certification time \(\le t_{\max}\) and memory \(\le m_{\max}\).


\item \textbf{Security boundaries.}\label{g:sec}
Disallow dynamic code injection in channel definitions; restrict to a whitelisted operator family. All MR artifacts execute in a sandbox with read-only corpora.


\item \textbf{Human oversight.}\label{g:human}
Escalate to human review when margins are within \(\delta\) of thresholds for \(N\) consecutive runs, or when novel prime signatures are introduced.
\end{enumerate}


\subsection*{Indexing \/ Cross-linking (Informative)}
Map \(U=X_\sigma + C_{\mathrm{lift}} + \Xi_{\mathrm{lift}}\) to the ACE-controlled form \(X + \sum_{p} w_p B_p\). Persist \((b_p,L_p)\) for certificate reuse and link entries to the PETC signature registry and ACE certificate store.


\subsection*{Verification Checklist (Complete per Release)}
\newcommand{\boxempty}{\(\square\)}
\begin{enumerate}[label=\boxempty\, C\arabic*., leftmargin=3.2em]
\item \textbf{GapLB:} computed and \(\gamma\ge\gamma_{\min}\). Artifacts attached.
\item \textbf{SlopeUB:} contraction margin \(\varepsilon>0\) certified. Runtime bound logged.
\item \textbf{Budgeting:} \(\sum_p |w_p|\le\tau\) and \(\|B_p\|\le b_p\), \(\mathrm{Lip}(B_p)\le L_p\) verified.
\item \textbf{PETC:} prime signatures validated; multiplicity/conservation checks passed.
\item \textbf{Ingest tests:} HS domain \/ boundedness \/ normality; multiplier; gap\/slope; lawfulness.
\item \textbf{Monitoring hooks:} telemetry for \((\gamma_{\min},\varepsilon,\tau)\) in place.
\item \textbf{Rollback:} golden set present; roll-forward\/back tested.
\item \textbf{Provenance:} hashes, versions, seeds, and dependency locks recorded.
\item \textbf{Security:} sandbox and whitelist active; no dynamic codepaths.
\item \textbf{Oversight:} thresholds \(\delta,N\) set; escalation route documented.
\end{enumerate}


\subsection*{Metadata Template (Copy into Release Notes)}
\begin{tabular}{@{}ll@{}}
\textbf{Artifact} & Meta‑Relativity Operator Stack (MR) \\
\textbf{DocID} & MR-IFMD-APPX-\textit{YYYYMMDD} \\
\textbf{Source Hash} & \texttt{<git-commit-or-file-hash>} \\
\textbf{Version} & \texttt{vX.Y.Z} \\
\textbf{ACE Budget} & \(\tau=\texttt{<value>}\) \\
\textbf{Gap Target} & \(\gamma_{\min}=\texttt{<value>}\) \\
\textbf{Contraction} & \(\varepsilon=\texttt{<value>}\) \\
\textbf{Bounds} & \(b_p=\texttt{<list>}\), \(L_p=\texttt{<list>}\) \\
\textbf{PETC Signatures} & \texttt{<registry-refs>} \\
\textbf{Reviewers} & \texttt{<names\/sign-offs>} \\
\textbf{Seeds} & \texttt{<rng-seeds>} \\
\textbf{Env} & \texttt{<platform\/deps lockfile>} \\
\end{tabular}


\subsection*{Operational Pseudocode (Informative)}
% This uses simple math notation to avoid extra packages; adapt to algorithm2e if desired.
\noindent\textbf{Pipeline:}
\begin{align*}
&\text{Input: MR artifact, budgets } (\tau, b_p, L_p),\; \gamma_{\min},\; \varepsilon.\\
&1.\; \text{Ingest tests }\Rightarrow \text{ pass else quarantine}.\\
&2.\; \text{Compute GapLB, SlopeUB under } (\tau, b_p, L_p).\\
&3.\; \text{If } \gamma\!<\!\gamma_{\min} \text{ or } \|U_{\mathrm{ACE}}\|\!>\!1-\varepsilon:\; \text{abort and rollback}.\\
&4.\; \text{Apply } U_{\mathrm{ACE}} = \Pi_{\mathrm{safe}} U \Pi_{\mathrm{safe}} \text{ with monitoring}.\\
&5.\; \text{Log certificates and telemetry; enable oversight triggers}.
\end{align*}


\subsection*{Change Control}
This section is versioned. Any modification requires a new DocID, reviewer sign-offs, and regeneration of certification artifacts. Deployments referencing this appendix must pin the exact DocID.

%========================
\section{Constitutional Foundations}
%========================

This work is grounded in a constitutional framework that specifies the fundamental architecture of physical law. The following axioms define the \emph{Meta-Relativity Principle} (MRP) and the \emph{Zeta--Multiplicity Axiom} (ZMA), establishing a foundation where mathematical structures are ontologically prior and arithmetic lawfulness constrains physical reality.

%------------------------
\subsection{The Meta-Relativity Principle (MRP)}
%------------------------

The \emph{Meta-Relativity Principle} (MRP) generalizes the principle of relativity by asserting that the laws of physics are invariant across different mathematical frames of reference. It posits that what we perceive as physical law is the invariant manifestation of a deeper, constitutional arithmetic structure.

\begin{description}
    \item[Axiom I (Mathematical Onticity).] The universe is isomorphic to a mathematical structure. The objects and relations of mathematics constitute the ontological substrate of physical reality; they are not merely descriptive models.

    \item[Axiom II (Frame Relativity).] Physical observables and laws are relative to a chosen mathematical frame $F$ (e.g., geometric, algebraic, categorical). Lawful transformations between frames preserve the constitutional invariants.

    \item[Axiom III (Prime Gating).] Prime numbers act as fundamental operators that gate access to lawful physical states. Admissible states must be expressible within a prime-indexed basis, ensuring a discrete, irreducible structure to lawfulness.

    \item[Axiom IV (Recursive Evolution).] A recursive operator $\Xi(t)$ governs dynamical evolution by driving systems toward stable attractors within the lawful subspace. Its action ensures coherence and stability against entropic decay over time.
\end{description}

\paragraph{Immediate Corollaries.}
These axioms yield several foundational consequences:
\begin{itemize}
    \item \textbf{Frame Covariance:} Physical predictions derived in one frame (e.g., the geometric AZ-TFTC formalism) must be isomorphic to those derived in another (e.g., the quantum information-theoretic Langlands Resonance framework).
    \item \textbf{Lawfulness as Prime Decomposability:} A physical state is lawful if and only if it admits a finite decomposition into a prime-indexed basis and remains stable under the action of $\Xi(t)$.
    \item \textbf{Constitutional Unification:} Experimental signatures (e.g., Casimir force spectra, cavity Q-factors) and arithmetic spectra (e.g., eigenvalues of the Zeta--Multiplicity Hamiltonian) are dual manifestations of the same constitutional invariants.
\end{itemize}

%------------------------
\subsection{The Zeta--Multiplicity Axiom (ZMA)}
%------------------------

The \emph{Zeta--Multiplicity Axiom} (ZMA) integrates the spectral theory of the Riemann zeta function directly into the constitutional framework, elevating its zeros to the status of fundamental physical actors.

\paragraph{Spectral Integration of Zeta Zeros.}
The non-trivial zeros $\rho = \frac{1}{2} + i\gamma$ of the Riemann zeta function constitute a spectral basis. A coupling operator $V_{p,\rho}$ maps these ``zeta modes'' to the prime-indexed basis states $\ket{p}$, intertwining the arithmetic and spectral landscapes.

\paragraph{Definition of the Lawful Subspace.}
The lawful subspace $\mathcal{H}_{\text{lawful}} \subset \mathcal{H}$ is defined as the image of a projector $P_{\text{CSL}}$ that enforces prime decomposability and stability under $\Xi(t)$:
\[
\mathcal{H}_{\text{lawful}} = \{ \ket{\psi} \, | \, \ket{\psi} = P_{\text{CSL}} \ket{\psi},\quad \ket{\psi} = \sum_p c_p \ket{p},\quad \Xi(t)\ket{\psi} = \lambda \ket{\psi} \}.
\]
This subspace is the arithmetic counterpart to the physically realizable state space of a resonant cavity with specific boundary conditions in the AZ-TFTC framework.

\paragraph{The abc-Gate as a Lawfulness Filter.}
The $abc$-conjecture provides a criterion for arithmetic balance. The operator $G_{abc}$ enforces this conjecture by projecting out states that violate it, acting as a necessary filter for lawfulness. This is formally analogous to how a fractal boundary in the AZ-TFTC experiment filters electromagnetic modes based on their overlap with the geometric potential $\Phi(\mathbf{x})$.

%------------------------
\subsection{Ethical and Epistemic Commitments}
%------------------------

The constitutional framework is not metaphysically neutral; it embeds specific ethical and epistemic commitments derived from the structure of lawfulness itself.

\begin{itemize}
    \item \textbf{Non-Arbitrariness:} The exclusion of unlawful states by $G_{abc}$ and $P_{\text{CSL}$ defines a form of cosmic non-arbitrariness. This constitutional prohibition against "unlawful drift" translates to an ethical imperative against coercive or exploitative actions, which are the social analogue of unstable, non-prime-decomposable states.

    \item \textbf{Epistemic Transparency:} The recursive verifiability of states under $\Xi(t)$ ensures that all lawful processes are transparent and justifiable, providing a foundation for trustworthy reasoning and prediction.

    \item \textbf{Multiplicity as Sanctity:} The preservation of prime-indexed diversity—the infinite variety of lawful states—is isomorphic to the principle of biodiversity and the sanctity of life. The constitutional requirement to preserve this multiplicity aligns physical law with a deep ethical commitment to sustainability and respect for complex systems.
\end{itemize}

The MRP and ZMA together form a coherent foundation that is simultaneously mathematically rigorous, physically testable, and ethically constraining. They demand a universe that is not only intelligible but also structured by a fundamental, lawful goodness.
\section{Field Equations and Operator Formulation}

This section presents the dynamical core of the Meta-Relativity framework: the operator equations that realize its constitutional axioms. We introduce two complementary formulations---the Zeta--Multiplicity Field Equation (ZMFE) and the Arithmetic Zeta--Topological Fractal Tensor Calculus (AZ-TFTC)---and demonstrate their physical equivalence.

\subsection{Zeta--Multiplicity Field Equation (ZMFE)}

\paragraph{State space.}
Let $\mathcal H_{\!P}:=\ell^2(\mathbb P)$ with orthonormal basis $\{\ket{p}\}_{p\in\mathbb P}$, and let
$\mathcal H_{\!\zeta}:=\ell^2(\mathcal Z)$ where $\mathcal Z$ indexes the nontrivial zeros $\rho=\tfrac12+i\gamma$.
Set the Hilbert space
\[
\mathcal H:=\mathcal H_{\!P}\oplus \mathcal H_{\!\zeta},\qquad
\langle (x,y),(x',y')\rangle := \langle x,x'\rangle_{\mathcal H_{\!P}}+\langle y,y'\rangle_{\mathcal H_{\!\zeta}}.
\]
We write $\ket{p}$ for prime basis vectors and $\ket{\rho}$ for zeta-zero basis vectors, with
$\{\ket{p}\}\cup\{\ket{\rho}\}$ orthonormal in $\mathcal H$.

The central object is the ZMFE Hamiltonian operator:
\begin{equation}\label{eq:ZM_Hamiltonian}
    H_{\mathrm{ZM}} = \Lambda_m I + \widehat{Z} + \widehat{M} + \left( \widehat{V}_{\mathrm{PR}} + \widehat{V}_{\mathrm{PR}}^\dagger \right) + \widehat{H}_{\mathrm{int}}(t),
\end{equation}
defined on a Hilbert space $\mathcal{H} = \operatorname{span}\{ \ket{p}, \ket{\rho} \}$ with orthonormal prime basis $\ket{p}$ and zeta basis $\ket{\rho}$ (indexed by the nontrivial zeros of the Riemann zeta function). Here:
\begin{itemize}
    \item $\widehat{Z} = \sum_\rho g(\rho) \ket{\rho}\bra{\rho}$ is the zeta spectral operator,
    \item $\widehat{M} = \sum_p m_p \ket{p}\bra{p}$ is the prime multiplicity operator,
    \item $\widehat{V}_{\mathrm{PR}}$ is the PIRTM coupling operator with matrix elements $v_{p,\rho} = \bra{p} \widehat{V}_{\mathrm{PR}} \ket{\rho}$,
    \item $\widehat{H}_{\mathrm{int}}(t)$ includes environmental couplings and penalty potentials that enforce constitutional constraints.
\end{itemize}
\paragraph{Standing assumptions.}
Throughout we assume $\widehat Z=\widehat Z^\ast$ (diagonal on $\{\ket{\rho}\}$ with real coefficients),
$\widehat M=\widehat M^\ast$ (diagonal on $\{\ket{p}\}$ with real coefficients), and $\widehat V_{PR}$ is bounded
with $\|\widehat V_{PR}\|<\infty$; $\widehat H_{\mathrm{int}}(t)$ is bounded and strongly measurable in $t$.
Under these hypotheses $H_{\mathrm{ZM}}$ is bounded self-adjoint and generates $e^{-itH_{\mathrm{ZM}}}$ on $\mathcal H$.

The dynamical evolution of a state $\ket{\Psi(t)}$ is then given by the Schrödinger equation:
\begin{equation}\label{eq:ZMFE}
    \left( H_{\mathrm{ZM}} + \Xi(t) \right) \ket{\Psi(t)} = i\hbar \, \partial_t \ket{\Psi(t)},
\end{equation}
subject to the constraint $P_{\mathrm{CSL}} \ket{\Psi(t)} = \ket{\Psi(t)}$, which projects the dynamics onto the lawful subspace $\mathcal{H}_{\mathrm{lawful}}$.

This formalism integrates arithmetic spectra directly into the generator of time evolution, extending foundational work on spectral operators in number-theoretic physics \cite{Connes1999, BostConnes1995}.

\subsection{AZ-TFTC Framework}

Complementary to the spectral ZMFE, the \emph{Arithmetic Zeta--Topological Fractal Tensor Calculus} (AZ-TFTC) provides a geometric formulation. It describes how arithmetic lawfulness manifests as an effective geometry for quantum fields.

The framework is built on a modified wave equation:
\begin{equation}\label{eq:AZTFTCEigenproblem}
    \left[ -c^2 \nabla^2 + V_{\mathrm{geo}}(\mathbf{x}) \right] \psi_n(\mathbf{x}) = \omega_n^2 \psi_n(\mathbf{x}),
\end{equation}
where the key innovation is the \emph{geometry potential}:
\begin{equation}\label{eq:GeoPotential}
    V_{\mathrm{geo}}(\mathbf{x}) = g \, \Phi(\mathbf{x}) + \eta \, \Phi(\mathbf{x}) R(\mathbf{x}) + \cdots.
\end{equation}
Here, $\Phi(\mathbf{x})$ is the \emph{geometric potency field}, a dimensionless scalar that encodes the local microstructure of spacetime (e.g., the fractal dimension of a boundary), and $R(\mathbf{x})$ is the Ricci scalar.

This potential parallels scalar-tensor theories of gravity and inflaton models in cosmology, but with a crucial difference: $\Phi(\mathbf{x})$ is not a fundamental field but an emergent quantity whose value is gated by prime resonance conditions. Fractal boundary conditions are used to model this confinement, drawing inspiration from both lattice gauge theory \cite{Wilson1974} and the physics of disordered systems \cite{Mandelbrot1982}.

The solutions to Eq.~\eqref{eq:AZTFTCEigenproblem} yield eigenmodes $\psi_n(\mathbf{x})$ that localize in regions of high geometric potency, producing a discrete spectrum of resonant frequencies $\omega_n$ that are arithmetic in origin.

\subsection{Equivalence and Unification}

A fundamental result of the Meta-Relativity framework is the operational equivalence between the spectral and geometric formulations:
\begin{equation}\label{eq:Equivalence}
    H_{\mathrm{ZM}} \quad \longleftrightarrow \quad H_{\mathrm{AZ-TFTC}} = -c^2 \nabla^2 + V_{\mathrm{geo}}.
\end{equation}
This duality implies that the prime-zeta couplings in $H_{\mathrm{ZM}}$ are isomorphic to the curvature-modulated $\Phi$-potentials in $H_{\mathrm{AZ-TFTC}}$.

This equivalence provides a powerful cross-verification mechanism. Physical predictions---particularly resonance structures in cavities with fractal boundaries, modifications to the Casimir effect, and spectral gaps---can be derived independently in both formalisms. This makes the framework highly testable. Experiments in superconducting microwave cavities and nano-optics are poised to detect these arithmetic-field predictions \cite{Jaffe2005}, potentially revealing the first direct signatures of number-theoretic lawfulness in quantum systems.

Together, the ZMFE and AZ-TFTC provide the first explicit dynamical instantiations of the Meta-Relativity Constitution, demonstrating how primes and zeta spectra yield coherent, testable physics.

\section{Mathematical Appendix: Proofs \& Bounds}

In this section, we provide technical results establishing the stability, convergence, and spectral bounds for the operators introduced in Sections~2 and~3. These results guarantee that the Meta-Relativity formalism is not only conceptually consistent but also mathematically well-posed.

\subsection{Operator Norm Bounds for $H_{\mathrm{ZM}}$}

\begin{proposition}
Let 
\begin{equation}
    H_{\mathrm{ZM}} = -\nabla^2 + \lambda \sum_{p \ \mathrm{prime}} \zeta_p(\hat{n},t)\, \Pi_p ,
\end{equation}
with $\Pi_p$ orthogonal projectors onto prime-gated subspaces of $L^2(\mathbb{R}^n)$.  
Then $H_{\mathrm{ZM}}$ is essentially self-adjoint on $C_c^\infty(\mathbb{R}^n)$, and satisfies the operator norm bound
\begin{equation}
    \| H_{\mathrm{ZM}} \| \leq \| \nabla^2 \| + |\lambda| \, C_\zeta ,
\end{equation}
where $C_\zeta = \sup_{p} \sup_{t} |\zeta_p(\hat{n},t)| < \infty$.
\end{proposition}

\begin{proof}
Self-adjointness follows from the Kato--Rellich theorem, since the prime-zeta term is relatively bounded with respect to the Laplacian. The operator norm estimate is immediate by bounding the zeta modes and applying triangle inequality over the prime projections.
\end{proof}

\subsection{Convergence of the Recursive $\Xi(t)$ Operator}

Define the recursive operator
\begin{equation}
    \Xi(t) = \sum_{k=0}^\infty \alpha_k(t) \, T^k ,
\end{equation}
where $T$ is a shift operator acting on prime-indexed multiplicity states, and $\alpha_k(t)$ are recursively defined coefficients.

\begin{theorem}
If $\sup_{t} \sum_{k=0}^\infty |\alpha_k(t)| < \infty$, then $\Xi(t)$ converges in the strong operator topology, and the family $\{ \Xi(t) \}$ is uniformly bounded:
\begin{equation}
    \| \Xi(t) \| \leq \sum_{k=0}^\infty |\alpha_k(t)| .
\end{equation}
\end{theorem}

\begin{proof}
The proof follows from the dominated convergence theorem applied to operator series on Banach spaces. Each term $T^k$ has unit norm, and hence the convergence reduces to absolute convergence of the coefficients $\alpha_k(t)$.
\end{proof}

\subsection{Existence and Uniqueness of the Lawful Subspace}

\begin{definition}
The \emph{lawful subspace} $H_{\mathrm{lawful}} \subset L^2(\mathbb{R}^n)$ is defined as the maximal subspace invariant under both $H_{\mathrm{ZM}}$ and $\Xi(t)$, subject to the prime decomposability constraint
\begin{equation}
    f \in H_{\mathrm{lawful}} \iff f = \sum_{p \ \mathrm{prime}} \Pi_p f .
\end{equation}
\end{definition}

\begin{theorem}
$H_{\mathrm{lawful}}$ exists uniquely as the intersection
\begin{equation}
    H_{\mathrm{lawful}} = \bigcap_{t} \mathrm{Dom}(\Xi(t)) \cap \bigcap_{p} \mathrm{Ran}(\Pi_p) .
\end{equation}
\end{theorem}

\begin{proof}
The prime projections $\{ \Pi_p \}$ form a commuting family of orthogonal projections, hence their intersection is a closed subspace. Since $\Xi(t)$ is uniformly bounded (Theorem 4.2), its domain is the entire Hilbert space. Thus the intersection is well-defined, closed, and invariant, ensuring uniqueness.
\end{proof}

\subsection{Error Estimates for Prime--Zeta Spectral Decomposition}

Let $\zeta(s)$ be approximated on the critical line $\Re(s) = 1/2$ by truncation at height $T$.  

\begin{proposition}
For $H_{\mathrm{ZM}}$ truncated to primes $p \leq P$ and zeta zeros $|\Im(s)| \leq T$, the spectral error satisfies
\begin{equation}
    \| H_{\mathrm{ZM}} - H_{\mathrm{ZM}}^{(P,T)} \| \leq C \left( P^{-\alpha} + T^{-\beta} \right) ,
\end{equation}
for some $\alpha, \beta > 0$, depending on the decay of prime sieve weights and explicit zero-density bounds \cite{Montgomery1973, IwaniecKowalski2004}.
\end{proposition}

\subsection{Casimir Multilayer Corrections}

Finally, we derive corrections to Casimir energies when arithmetic layers are imposed between conducting plates. For separation $d$ with $N$ prime-indexed boundary layers, the corrected Casimir energy is
\begin{equation}
    E_{\mathrm{Casimir}}(d,N) = - \frac{\pi^2}{720 d^3} \left[ 1 + \sum_{p \leq p_N} \frac{\gamma_p}{p^2} \right] ,
\end{equation}
where $\gamma_p$ are multiplicity coupling coefficients associated with each prime gate.  

\begin{proof}[Sketch of Proof]
Start from the standard mode sum for Casimir energy. Impose prime-gated boundary conditions, leading to modified eigenvalue spacing proportional to $1/p$. The correction term then arises as a multiplicative series expansion, with convergence ensured by $\sum_p p^{-2} < \infty$.
\end{proof}

\section{Experimental Predictions}

The validity of the Meta-Relativity framework hinges on its empirical falsifiability. This section delineates concrete, testable predictions derived from both the AZ-TFTC and ZMFE formulations. These predictions are designed to be probed using established techniques in cavity quantum electrodynamics (QED), Casimir force metrology, and nanoscale resonator engineering.

\subsection{AZ-TFTC Predictions (P1--P5)}

The Algebraic Zeta–Topological Fractal Tensor Calculus (AZ-TFTC) yields the following primary experimental signatures:

\begin{description}
    \item[P1: Prime-Resonant Mode Stabilization.] 
    Microwave and optical cavities with boundaries engineered to have a high geometric potency $\Phi$ (e.g., fractal patterns with dimension $D$ tuned to a prime-related value) will exhibit a \emph{mode-dependent enhancement of the quality factor $Q$}. Modes with a high spatial overlap $\langle \psi_n | \Phi | \psi_n \rangle$ with the fractal region will show a measurably higher $Q$ than modes in a smooth cavity or modes with low overlap. This is a direct signature of prime-gated lawfulness suppressing decoherence.

    \item[P2: Casimir Force Deviations in Multilayer Systems.]
    The Casimir force between surfaces coated with prime-structured multilayer stacks will deviate from the standard Lifshitz theory prediction. The fractional deviation $\Delta F_C / F_C$ is predicted to scale with a sum over primes constitutive of the stack geometry:
    \[
    \frac{\Delta F_C}{F_C} \propto \sum_{p \leq p_{\mathrm{max}}} \frac{\gamma_p}{p^2},
    \]
    where $p_{\mathrm{max}}$ is a cutoff prime related to the smallest layer thickness. This provides a clear arithmetic signature distinguishable from conventional material or roughness effects.

    \item[P3: Spectral Gaps at Prime-Logarithmic Intervals.]
    The resonant frequency spectrum $\{\omega_n\}$ of a prime-optimized cavity will exhibit anomalous gaps or suppressions at frequencies corresponding to $\omega \sim \omega_0 \log p$, where $p$ is a prime number. These non-geometric spectral features are a direct consequence of the $abc$-gate filtering unlawful fluctuations.

    \item[P4: Anomalous Frequency Shifts.]
    The resonant frequencies of cavity modes will experience shifts $\delta\omega_n$ proportional to their overlap with the geometric potency field:
    \[
    \frac{\delta\omega_n}{\omega_n} = -\frac{g}{2\omega_n^2} \langle \psi_n | \Phi | \psi_n \rangle.
    \]
    This shift is mode-specific and provides a quantifiable measure of the coupling strength $g$.

    \item[P5: Vacuum Energy Density Modification.]
    The local density of states (LDOS) $g(\omega, \mathbf{x})$ near a fractal boundary will be altered, leading to a measurable shift in the Casimir energy density and the Purcell effect for emitters placed near such a surface. This represents a direct manipulation of vacuum energy via arithmetic geometry.
\end{description}

\subsection{ZMFE Experimental Consequences}

The Zeta–Multiplicity Field Equation (ZMFE) provides a spectral counterpart to the AZ-TFTC's geometric predictions, leading to convergent experimental consequences:

\begin{itemize}
    \item \textbf{Prime–Mode Spacing:}  
    The eigenvalue spectrum of the ZMFE Hamiltonian $H_{\mathrm{ZM}}$ clusters at values corresponding to $\lambda_n \sim \lambda_0 \log(p)$. The physical manifestation is a corresponding clustering of cavity resonance frequencies at $\omega_n \sim \omega_0 \log(p)$, providing an alternative basis for predicting P3.

    \item \textbf{Lawful Resonance Bands:}  
    After applying the $G_{abc}$ filter, the ZMFE spectrum decomposes into distinct, stable ``lawful bands.'' The physical analogue is the experimental observation of a subset of cavity modes (those with high prime-overlap) exhibiting anomalously narrow linewidths and high stability (as in P1).

    \item \textbf{Cross-Verification:}  
    The ZMFE and AZ-TFTC formulations must yield consistent predictions for the same physical observable (e.g., the $Q$-factor of the $\mathrm{TE}_{211}$ mode). This consistency across independent mathematical frameworks is a critical test of the entire Meta-Relativity constitution.
\end{itemize}

\subsection{Verification Plan: Casimir Force and Cavity QED}

The most direct path to verification involves a combined study of Casimir forces and resonant cavity spectra.

\paragraph{Casimir Force Metrology.}
\begin{enumerate}
    \item \textbf{Fabrication:} Construct parallel-plate or sphere-plate Casimir apparatuses where one surface is coated with a multilayer stack or a fractal pattern designed with prime-indexed feature sizes.
    \item \textbf{Measurement:} Perform high-precision force measurements across separation distances $d$ from 50 nm to 1 $\mu$m, achieving sub-percent accuracy to resolve the predicted deviations $\Delta F_C$.
    \item \textbf{Analysis:} Fit the residual force, after subtracting the standard Lifshitz prediction, to the arithmetic correction form $\sum_p \gamma_p / p^2$. A statistically significant fit would constitute a discovery.
\end{enumerate}

\paragraph{Cavity Resonance Spectroscopy.}
\begin{enumerate}
    \item \textbf{Fabrication:} Machine high-$Q$ microwave and optical cavities with identical gross geometry but different inner surface treatments: one smooth (control) and one with a prime-optimized fractal etch.
    \item \textbf{Measurement:} Use vector network analysis and resonance fluorescence techniques to precisely measure the spectrum $\{\omega_n\}$ and quality factors $\{Q_n\}$ for all modes in a given frequency band.
    \item \textbf{Analysis:} Identify modes with anomalous properties: (i) those with frequency shifts $\delta\omega_n$ consistent with P4, (ii) those with enhanced $Q$ consistent with P1, and (iii) those whose frequencies align with prime-logarithmic spacing consistent with P3.
\end{enumerate}

\paragraph{Implications of Verification.}
A positive result would demonstrate that arithmetic lawfulness is a tangible physical phenomenon, opening the field of \emph{vacuum energy engineering}. This could enable new technologies based on controlling decoherence and vacuum fluctuations through geometric design, with potential applications in quantum computing, ultra-stable oscillators, and fundamental metrology.
\section{Interdisciplinary Integration}

The constitutional framework presented here does not stand in isolation. 
It actively converges with independent lines of research across topology, cosmology, number theory, and quantum information. 
These integrations strengthen the thesis that primes and multiplicity are not merely abstract constructs but universal regulators of lawfulness across physical, mathematical, and computational domains. 

\subsection{CronNet-Holo Initiative Contributions}

The CronNet-Holo annex \cite{CronNetAnnex} has outlined a programmatic connection between deep topological invariants and the structure of time itself. 
Three strands are particularly relevant to the present work: 

\begin{enumerate}
    \item \textbf{Conway Knot Topology:} The Conway knot, as a minimal nontrivial example of topological entanglement, provides a candidate model for the ``irreducible torsion'' in temporal recursion. This resonates with the $\Xi(t)$ operator formalism of the Meta-Relativity Principle, where recursion embodies the lawful generation of time-frames. 
    \item \textbf{Dark Flow Alignment:} Large-scale anisotropies, such as the hypothesized dark flow, suggest preferred frames or boundary alignments in the cosmic vacuum. Within AZ-TFTC, these can be interpreted as large-scale fractal constraints on cavity-like domains. 
    \item \textbf{10-adic Discretization:} The use of $10$-adic sieving to discretize temporal flows parallels the prime sieving inherent in the ZMFE Hamiltonian. Both enforce granular, prime-gated filters on otherwise continuous spectra.
\end{enumerate}

Thus, CronNet-Holo provides an independent but consistent set of constructs that reinforce the constitutional axioms of MRP and ZMA. 

\subsection{Langlands Correspondence and Automorphy}

The Langlands program has long been seen as a deep unifying framework linking number theory and representation theory. 
In our context, it serves as a regulator and stabilizer of spectral behavior. 

\begin{itemize}
    \item \textbf{Automorphic Forms:} As natural eigenfunctions of the Laplacian on arithmetic domains, automorphic forms act as spectral regulators that map directly onto the ``lawful resonance bands'' predicted by the ZMFE. 
    \item \textbf{Langlands Prism:} Prior work has proposed the Langlands prism as a categorical structure within the Ξ-Constitution (node 137). This prism functions as a higher-order filter, enforcing automorphic symmetries on the lawful subspace. 
\end{itemize}

Recent developments in Langlands-encoded resonance \cite{LanglandsResonance} suggest that these same symmetries can be physically instantiated in quantum error-correcting codes and gravitational spectra, further corroborating the isomorphic mapping between AZ-TFTC and LRC.

\subsection{Quantum Information \& Computation Links}

A key corollary of the constitutional framework is that lawfulness must be verifiable not only mathematically but also computationally. 
This naturally points toward quantum information science as both a proving ground and an application domain. 

\begin{itemize}
    \item \textbf{zk-SNARK Lawful Verification:} Zero-knowledge succinct non-interactive arguments of knowledge (zk-SNARKs) provide cryptographic protocols for enforcing lawfulness without disclosure. In the context of ZMFE, they correspond to verifying that a given eigenstate belongs to $\mathcal{H}_{\mathrm{lawful}}$ without requiring full disclosure of the prime decomposition. 
    \item \textbf{Prime-Indexed Circuits:} The recursion operator $\Xi(t)$ and the abc-gate naturally induce circuits indexed by primes, which can be implemented in arithmetic-friendly programming frameworks such as Circom or Noir. These circuits embody lawful computation by construction, ensuring that any output is consistent with the constitutional axioms. 
\end{itemize}

This integration closes the loop: from physical cavity experiments (AZ-TFTC), to arithmetic operators (ZMFE), to topological time (CronNet-Holo), to automorphic number theory (Langlands), and finally to quantum information and lawful computation. 
Each field provides a distinct ``frame'' of the same constitutional substrate, in full alignment with the Meta-Relativity Principle. 

\section{Summary and Outlook: The Multiplicity Relativity Program}

\subsection{Constitutional Developments}

The Multiplicity Relativity Program (MRP) has advanced from a speculative framework into a constitutionally rigorous and experimentally falsifiable research paradigm. Its foundation rests on the recognition that lawful physical states are confined to a prime-gated, multiplicative subspace
\[
\mathcal{H}_{\text{lawful}} \subset \ell^2(\mathbb{N}),
\]
spanned by multiplicative characters and stabilized by a recursive contraction operator $\Xi(t)$ \cite{Edwards1974,Connes1999}. This insight yields a new set of axioms extending traditional physics:

\begin{itemize}
    \item \textbf{Axiom M1 (Prime-Gated State Space):} Physical states are lawful iff they belong to $\mathcal{H}_{\text{lawful}}$.
    \item \textbf{Axiom M2 (Recursive Evolution):} Evolution contracts toward $\mathcal{H}_{\text{lawful}}$ under $\Xi(t) = e^{tG}$, with $G = -\sum_p \beta_p (I - P_p)$.
    \item \textbf{Axiom M3 (Frame Relativity):} Observables are projections of lawful states into isomorphic frames (geometric, number-theoretic, quantum-informational).
    \item \textbf{Axiom M5 (Emergent Causality):} Stable macroscopic plateaus correspond to recursive fixed points of $\Xi(t)$.
    \item \textbf{Axiom M6 (Information-Form Equivalence):} Physical constants (e.g. mass gaps) are frame-dependent and may drift under prime resonance.
\end{itemize}

\subsection{Mathematical Core and Theorems}

The mathematical kernel of the program is Prime-Indexed Recursive Tensor Mathematics (PIRTM) \cite{CitizenGardens:2023a}, which introduces the universal multiplicity constant $\Lambda_m$ as a coupling invariant across frames \cite{CitizenGardens:2023c}. 

\begin{theorem}[Lawful Stabilization]
Let $f \in \ell^2(\mathbb{N})$. Then $\lim_{t \to \infty} \Xi(t) f \in \mathcal{H}_{\text{lawful}}$. Moreover, $\Xi(t)$ is a bounded, contractive semigroup.
\end{theorem}

This theorem establishes the closure and stability of lawful states, providing the bridge between number theory and physical instantiation.

\subsection{Experimental Program and Pre-Registered Tests}

Three decisive experimental predictions connect MRP to observation:

\begin{enumerate}
    \item \textbf{P1 (Neutron $\beta$-decay endpoint shift):} Tests H$_{\text{strong}}$---the hypothesis of an embedded positron or charge cohomology label in the neutron. A measurable shift $B$ in the endpoint energy would constitute a violation of Standard Model predictions \cite{Weinberg:1995mt,ParticleDataGroup:2022pth,Wietfeldt:2011ue}.
    
    \item \textbf{P2 (Prime cavity stability):} Tests H$_{\text{weak}}$ by measuring anomalous Q-factor uplifts in prime-indexed microwave modes. Contractivity of $\Xi(t)$ predicts reduced dissipation for primes, analyzed under the Benjamini–Hochberg procedure \cite{Haroche:2013bs,Benjamini:1995}.
    
    \item \textbf{P3 (He-BEC coherence plateaus):} Tests macroscopic frame-relativity by identifying drive-response plateaus in superfluid helium, analogous to recursive fixed points \cite{Day:2009zz}.
\end{enumerate}

All protocols are pre-registered with explicit statistical decision rules: $5\sigma$ significance for P1, FDR-controlled detection for P2, and reproducible plateau thresholds for P3.

\subsection{Results to Date}

Preliminary simulations indicate that the contractive flow $\Xi(t)$ suppresses noise in synthetic prime-indexed cavities at the predicted 5--10\% level. Initial design studies for the neutron endpoint experiment confirm feasibility of sub-keV calibration precision. Full experimental data are pending collaboration with ultracold neutron groups and cavity QED laboratories.

\subsection{Implications and Future Directions}

The program yields three possible outcomes:

\begin{itemize}
    \item \textbf{P1 positive:} A paradigm-shifting revision of neutron structure, requiring extension beyond the Standard Model, in resonance with Weinberg’s field theory foundation \cite{Weinberg:1995mt}.
    \item \textbf{P2/P3 positive with P1 null:} Validation of the lawful subspace framework without requiring embedded positrons, supporting a reformulation of quantum coherence via multiplicative recursion \cite{CitizenGardens:2023b}.
    \item \textbf{All null:} Multiplicity Relativity is falsified at current sensitivity, and its constitutional axioms require revision.
\end{itemize}

Speculative appendices connect this framework to Langlands duality \cite{Langlands:1967}, the Riemann hypothesis \cite{Berry1979}, macroscopic quantum coherence in biology \cite{Hameroff:2014}, and Casimir experiments probing multiplicative vacua \cite{Lamoreaux:1996wh,Mohideen1998}.

\subsection{Conclusion}

The Multiplicity Relativity Program provides a falsifiable, constitutionally rigorous paradigm linking prime multiplicativity to physics. Its success or failure hinges on imminent experimental results. Either outcome will significantly advance the discourse on the lawful foundations of reality.

\section*{Appendix S (Revised): Sommerfeld--Meta-Relativity Operator \& Testable Prediction}

\subsection*{S.1 Lawful Sector, Prime Register, and Zeta Kernel}

\paragraph{Hilbert space.}
We work on
\[
\mathcal H \equiv L^2(\mathbb R^3;\mathbb C^2)\ \widehat{\otimes}\ \ell^2(\mathbb P),
\quad
\Psi(x)=\sum_{p\in\mathbb P}\psi_p(x)\,|p\rangle.
\]

\begin{definition}[Zeta kernel $X(\sigma)$]
For $\sigma>1$, define on the prime register
\[
(\mathsf Z_\sigma f)(p)\ :=\ \sum_{q\in\mathbb P}\frac{(\log p\,\log q)^{1/2}}{(pq)^{\sigma/2}}\,f(q),\qquad
X(\sigma):=I_{L^2}\otimes \mathsf Z_\sigma.
\]
Then
\[
\|\mathsf Z_\sigma\|_{HS}^2=\sum_{p,q}\frac{\log p\,\log q}{(pq)^{\sigma}}
=\Big(\sum_{p}\frac{\log p}{p^\sigma}\Big)^2<\infty \ \ (\sigma>1),
\]
so $\mathsf Z_\sigma$ (hence $X(\sigma)$) is Hilbert--Schmidt, positive, compact, and self-adjoint.
\end{definition}

\begin{remark}[Why primes?]
Primes are the atoms of multiplicative structure; $\mathsf Z_\sigma$ is the Gram operator of the family $p^{-\sigma/2}$ with log-weights, giving a multiplicative-factorization--aware spectral gate. Lawfulness $\Leftrightarrow$ controlled multiplicative composability.
\end{remark}

\subsection*{S.2 Lawful Sectors (two compatible options)}

\paragraph{Option A (σ-independent, safe).}
Fix $\sigma>1$ and $\epsilon>0$. Define
\[
\mathcal H_{\mathrm{lawful}}
:=\Big\{\Psi:\ \sup_{x\in\mathbb R^3}\ \sum_{p}\frac{\|\psi_p(x)\|^2}{(\log p)^{2+\epsilon}}<\infty\Big\}.
\]
Then $X(\sigma)$ and the gates below are bounded on $\mathcal H_{\mathrm{lawful}}$; truncations in $P$ converge in norm.

\paragraph{Option B (σ-weighted, allows }\tfrac12<\sigma\le 1\text{).}
For $\tfrac12<\sigma\le 1$ set weights $w_p:=\dfrac{p^\sigma}{\log p}$ and
\[
\mathcal H_{\mathrm{lawful}}^{(\sigma)}
:=\Big\{\Psi:\ \sup_{x}\sum_{p}\frac{\|\psi_p(x)\|^2}{w_p}<\infty\Big\}.
\]
Then the matrix $k_{pq}=\dfrac{\sqrt{\log p\,\log q}}{(pq)^{\sigma/2}}$ defines a bounded form by Schur/Cauchy–Schwarz with $w_p$ (see Lemma S.\ref{lem:schur}).

\subsection*{S.3 Sommerfeld Gate and SMR Operator}

Let $H_D$ be the Dirac–Coulomb Hamiltonian; $\mathcal M_{j\ell}$ the spin–orbit projector.

\begin{definition}[Sommerfeld gate $\Delta_\alpha$]
For cutoff $P$ and geometry potency $\Phi\in(0,1)$, set
\[
(\Delta_\alpha \Psi)(x)
=\alpha\,\Phi\sum_{p,q\le P}\frac{(\log p\,\log q)^{1/2}}{(pq)^{\sigma/2}}\,
\mathcal M_{j\ell}\,\psi_q(x)\,|p\rangle.
\]
\end{definition}

\begin{definition}[Recursive stabilizer and SMR operator]
Let $\Xi(t)=e^{-t\mathsf R}$ be a contraction semigroup on $\ell^2(\mathbb P)$ with $\mathsf R\ge 0$ trace-class. Define
\[
H_{\mathrm{SMR}}(\sigma,t):=\big(H_D\otimes I + X(\sigma)+\Delta_\alpha+\Xi(t)\big)\big|_{\mathcal H_{\mathrm{lawful}}}
\quad\text{(Option A)},
\]
or with $\mathcal H_{\mathrm{lawful}}^{(\sigma)}$ under Option B.
\end{definition}

\begin{theorem}[Self-adjointness and unitary flow]\label{thm:sa-new}
Let $\sigma>1$ (Option A), or $\tfrac12<\sigma\le 1$ with $\mathcal H_{\mathrm{lawful}}^{(\sigma)}$ (Option B).
Then $H_{\mathrm{SMR}}(\sigma,t)$ is essentially self-adjoint on the Dirac core and generates a strongly continuous unitary group.
\emph{Sketch.} $H_D$ is essentially self-adjoint on the standard core. $X(\sigma)$ is compact for $\sigma>1$ (HS); $\Delta_\alpha$ and $\Xi(t)$ are bounded; under Option B boundedness follows from Lemma~\ref{lem:schur}. Apply Kato–Rellich and Stone.
\end{theorem}

\subsection*{S.4 Derived Fine Structure (unchanged structure, clarified gain)}

For hydrogenic $\ket{n\kappa m}$ with $j=|\kappa|-1/2$,
\[
\Delta E_{n\kappa}^{(1)}=\alpha\,\Phi\,S(P,\sigma)\,C_{n\kappa},\qquad
S(P,\sigma):=\sum_{p\le P}\frac{\log p}{p^\sigma},
\]
with
\[
C_{n\kappa}=\frac{1}{n^3}\!\left(\frac{1}{j+\tfrac12}-\frac{3}{4n}\right)+\mathcal O(n^{-4}).
\]
Hence
\[
\boxed{\ \Delta E_{n,j}^{(1)}=\alpha\,\Phi\,S(P,\sigma)\,\frac{1}{n^3}
\!\left(\frac{1}{j+\tfrac12}-\frac{3}{4n}\right)+\mathcal O(\alpha^2\Phi^2)\ }.
\]

\begin{remark}[Finite-window gain vs convergence]
For $\sigma\le 1$, $S(P,\sigma)$ grows with $P$; here $P$ is an \emph{experimental window} and is never sent to $\infty$—the growth is a tunable gain. For $\sigma>1$, $S(P,\sigma)\to S(\infty,\sigma)$ as $P\to\infty$ and may be treated as a calibration constant.
\end{remark}

\begin{remark}[Sign]
$\Delta E$ follows $\mathrm{sgn}(\Phi)$ since the kernel is positive. Geometries with effective negative $\Phi$ flip the sign of the shift.
\end{remark}

\subsection*{S.5 ACE-Style Spectral Certification}

Let $\lambda_{\min}(\sigma;P)$ be the smallest nonzero eigenvalue of $\mathsf Z_\sigma$ on the chosen window $\{p\le P\}$. Define the perturbation scale
\[
\eta_{\mathrm{MR}}:=\alpha\,\Phi\,\|\mathsf Z_\sigma\|\,\|\mathcal M_{j\ell}\|.
\]
We certify first-order validity by
\[
\eta_{\mathrm{MR}}\ \le\ \varepsilon\,\lambda_{\min}(\sigma;P),\qquad \varepsilon\lesssim 0.1.
\]
Report $(P,\sigma,\lambda_{\min},\|\mathsf Z_\sigma\|,\alpha\Phi)$ with each dataset.

\subsection*{S.6 Numerical Values for $S(P,\sigma)$ (illustrative)}

\begin{table}[h]
\centering
\caption{Representative $S(P,\sigma)=\sum_{p\le P}\log p/p^\sigma$. Values increase with $P$ for $\sigma\le 1$ (finite-window gain).}
\label{tab:Svals-new}
\begin{tabular}{@{}lrrrrr@{}}
\toprule
$\sigma\backslash P$ & 100 & 300 & 1000 & 3000 & 10000 \\
\midrule
0.90 & 9.261 & 15.701 & 26.885 & 41.842 & 63.917 \\
1.00 & 8.360 & 13.983 & 23.506 & 36.275 & 55.186 \\
1.10 & 7.621 & 12.529 & 20.610 & 31.036 & 45.333 \\
\bottomrule
\end{tabular}
\end{table}

\subsection*{S.7 Experimental Notes (calibration and falsification)}

\paragraph{\(\kappa_{\mathrm{FS}}\) calibration (one-time).}
Determine $\kappa_{\mathrm{FS}}$ by measuring the $2p$ doublet in a smooth cavity (baseline) and with the fractal insert at a low-gain setting (small $\Phi S$). Fit the observed $\Delta(\Delta\nu)$ to
\[
\frac{\Delta(\Delta \nu)}{\Delta \nu_{\mathrm{Dirac}}}
=\kappa_{\mathrm{FS}}\ \alpha\,\Phi\,\frac{S(P,\sigma)}{n^3}
\quad(n=2),
\]
then hold $\kappa_{\mathrm{FS}}$ fixed for all prime-window scans.

\paragraph{Prime-window falsification.}
Besides randomizing mode order, perform a \emph{scrambled-weights} run (same amplitudes, weights shuffled across primes). The prime-correlated signature must disappear.

\subsection*{S.8 Proof Fragments}

\begin{prop}[HS threshold]\label{prop:hs}
$\mathsf Z_\sigma$ is Hilbert--Schmidt iff $\sigma>1$.
\emph{Proof.}
$\|\mathsf Z_\sigma\|_{HS}^2=(\sum_p \log p/p^\sigma)^2$; the scalar series converges iff $\sigma>1$ by comparison with $\sum_p 1/p^{\sigma-\delta}$ and the PNT estimate $\pi(x)\sim x/\log x$.
\qed
\end{prop}

\begin{lemma}[Schur test for $\tfrac12<\sigma\le 1$]\label{lem:schur}
With $w_p=p^\sigma/\log p$, the kernel $k_{pq}=\sqrt{\log p\,\log q}/(pq)^{\sigma/2}$ defines a bounded operator on $\ell^2(w^{-1})$.
\emph{Proof.}
For $f\in\ell^2(w^{-1})$, by Cauchy–Schwarz
\[
\sum_q |k_{pq}|\,|f_q|
\le \Big(\sum_q k_{pq}^2 w_q\Big)^{1/2}\Big(\sum_q \frac{|f_q|^2}{w_q}\Big)^{1/2}.
\]
But $k_{pq}^2 w_q = \dfrac{\log p\,\log q}{(pq)^\sigma}\cdot\dfrac{q^\sigma}{\log q}=\dfrac{\log p}{p^\sigma}$,
so $\sum_q k_{pq}^2 w_q=\dfrac{\log p}{p^\sigma}$ and $\sum_p \dfrac{\log p}{p^\sigma}<\infty$ for any fixed $P$ (or is controlled uniformly on windows). Symmetric bound in $p$ gives the Schur criterion.
\qed
\end{lemma}

\bibliographystyle{unsrt}
\bibliography{references}
\clearpage

\end{document}
